\documentclass[a4paper,10pt]{article}
\usepackage[utf8]{inputenc}
\usepackage[TS1,T2A]{fontenc}
\usepackage[english,russian]{babel}

%opening
\title{Кореляционный метод детектирования на основе параметрического моделирования опорного сигнала}
\author{Ушаков Виталий}
\date{2025-01-08}

\begin{document}

\maketitle
\begin{abstract}
Метод детектирования сигнала по пороговому уровню дает много ложных результатов при значения соотношения сигнал/шум (SNR) менее 2. А при соотношениях близких к 1 и мене даже ``вручную'', визуально определить время прихода сигнала становится затруднительно или вообще невозможно.

На практике такие зашумленные записи встречаются. В таком случае встает необходимость задействовать более мощные источники возбуждения, что ведёт к дополнительным затратам.

Предложенный метод позволяет детектировать полезный сигнал с приемлимой точностью даже при значениях SNR менее 2.
\end{abstract}

\newpage
\section{Описание метода}
По записи с приёмника, ближайшего к источноку возбуждения, построить параметрическую модель опорного сигнала, который не будет иметь шумовую составляющую. Провести корреляцию всей записи с этим опорным сигналом для детектирования искомых моментов времени.

\newpage
\section{Испытание методом моделирования}
Сгенерировать шумовой сигнал с заданными величинами пиковой или среднеквадратичной амплитуды, состоящий из гармоник, входящих в рабочий диапазон частот приемников (от 5 до 120 Гц).
Провести испытания при различных соотношениях Сигнал/Шум. Сперва по соотношению пиковых значений амплитуд (S/Np-p).

Количество вариаций шума при тех же параметрах 1000.
Результат - процент отклонений от точного местоположения.


\subsection{Эксперимент 1.}

S/Npp = 2;

ref: A=1, len=0.5, fr=20.0, att=20.0

\newpage
\section{Моделирование опорного сигнала}

Опорный сигнал можно представить в виде функции
\[ S_{ref}(t) = A\sin{\omega t}\cdot e^{-\alpha t} \]

$\omega$ - угловая частота, равная $2\pi f$

$\alpha$ - коэффициент затухания

$A$ - амплитуда


Определять их будем по записи с приемника, расположенного вблизи источника возбуждения. Запись хоть и имеет большое значение SNR, тем не менее содержит шумы. Следовательно при вычислении характеристик будут погрешности. На данном этапе важно определить допустимый порядок величин этих погрешностей.

\subsection{Определение допустимых отклонений параметров модели}
Определим как изменяется точность детектирования при отклонении параметров опорного сигнала $S_{refvar}$ от исходного $S_0$. Для этого зададим неизменяемый $S_0$, параметры же $S_{refvar}$ будем менять, наблюдая результат их взаимной корреляции $S_{cor}$. Очевидно, что отклонение будет миниимальным при одинаковых параметрах сигнала. За допустимое примем отклонение не более одного переиода дискретизации.

TODO: ТЕСТЫ


\subsection{Определение частоты}
Частота обратна периоду. Период равен 2-м полупериодам. На каждом полупериоде сигнал достигает локального максимума. По ним и будем определять частоту. Простой пиковый детектор часто ошибается, потому моменты времени будем определять по наибольшим локальным значениям площади сигнала в пределах заданного окна.

\subsection{Вычисление коэффициента затухания}

Сигнал затухает экспоненциально. Это можно описать уравнением
\[f(t) = A_0e^{-\alpha t}\]
где $A_0$ - начальная амплитуда,
$\alpha$ - коэффициент затухания.

Значения функции в определенные моменты времени $t_1$ и $t_2$ нам известны из предыдущих вычислений. Коэффициент определим из соотношения:
\[ \frac{A_0e^{-\alpha t_2}}{A_0e^{-\alpha t_1}} = \frac{f(t_2)}{f(t_1)} \]
Тогда:
\[ \alpha = -\frac{\ln|\frac{f(t_2)}{f(t_1)}|}{t_2 - t_1} \]

\subsection{Вычисление начальной амплитуды}

Подставив известные значения $t_1$ и $\alpha$ в формулу
$f(t_1) = A_0e^{-\alpha t_1}$, определим начальную амплитуду:
\[ A= \frac{f(t_1)}{e^{-\alpha t_1}}\]

\end{document}
